\section{Motivation}

If there is one mathematical habit that defines modern computing, it is this: \emph{turn your data into a
vector}. A greyscale image of size $1000\times 1000$ is a list of a million pixel brightnesses — a vector
in $\R^{1\,000\,000}$. An audio clip is a vector of samples. A user of a streaming service is a vector of
ratings. A word, inside a modern language model, is a vector of a few hundred coordinates learned from
text. Once data is a vector, geometry becomes a computational tool: similar users are \emph{nearby}
vectors, recommendation is finding neighbours, and classification is separating regions of a space that we
cannot draw but can compute in.

The operations we axiomatise in this chapter are meaningful operations on data. Adding vectors and scaling
them is how images are blended and audio tracks are mixed. More strikingly, in word-vector models the sum
and difference of vectors capture \emph{meaning}: the famous relation
\[
\text{\itshape king} - \text{\itshape man} + \text{\itshape woman} \approx \text{\itshape queen}
\]
is literally vector arithmetic in the sense of this chapter.

The deeper notions of the chapter answer a question every computer scientist should ask about data:
\emph{how many numbers does it really take to describe it?}
\begin{itemize}
\item A \emph{spanning set} is a collection of building blocks capable of expressing everything in sight.
\item \emph{Linear independence} means no building block is redundant: none can be reconstructed from the
others, so none can be dropped for free.
\item A \emph{basis} is a minimal, lossless set of building blocks, and the \emph{dimension} — the size of
any basis — is the honest number of degrees of freedom of the data.
\end{itemize}
Real data rarely uses all of its apparent dimensions: photographs of faces occupy a tiny sliver of
$\R^{1\,000\,000}$, and that sliver is well approximated by low-dimensional \emph{subspaces}. This single
observation is why image compression, denoising and recommender systems work at all, and we will be able
to \emph{compute} such dimensions once matrix rank is available in Chapter~\ref{ch:matrices}.

One more remark on abstraction. We will not define vectors as ``arrows'' or as ``tuples of numbers'', but
as elements of \emph{any} set equipped with a well-behaved addition and scaling — programming against an
interface again, as in the previous chapter. The payoff is immediate: tuples, polynomials, matrices and
functions all satisfy the axioms, so every theorem proved below holds for all of them at once.

\section{Vector spaces}

You are already familiar with the two-dimensional vector space
\[
\R^2 = \{(x,y)\mid x,y\in \R\}.
\]
In $\R^2$ vectors can be added and multiplied by a number called a \emph{scalar}:
\[
(x_1,y_1)\oplus(x_2,y_2) = (x_1+x_2,\, y_1+y_2), \qquad
\lambda \odot (x,y) = (\lambda x,\, \lambda y) \text{ for } \lambda\in \R.
\]
The abstract notion of a vector space axiomatises exactly these two operations.

\begin{definition}
Let $\K$ be a field. A set $\V$ together with two operations:
\begin{itemize}
\item \emph{vector addition} $+\colon \V\times \V\to \V$ and
\item \emph{scalar multiplication} $\cdot\colon \K\times \V\to \V$
\end{itemize}
is called a \emph{vector space} over the field $\K$ if it satisfies the following conditions:
\begin{itemize}
\item for any $\uu,\vv,\ww\in \V$ we have $\uu+\vv=\vv+\uu$ and $\uu+(\vv+\ww) = (\uu+\vv)+\ww$;
\item there is $\0\in \V$, called the \emph{zero vector}, such that for any $\vv\in\V$ we have
$\vv+\0 =\0+\vv=\vv$;
\item for any $\vv\in\V$ there is $-\vv\in\V$ such that $\vv+(-\vv) = (-\vv)+\vv = \0$;
\item for any $\uu,\vv\in \V$ and any scalar $\lambda\in \K$ we have
$\lambda\cdot(\uu+\vv)=\lambda\cdot \uu + \lambda\cdot\vv$;
\item for any $\uu\in \V$ and any scalars $\lambda_1,\lambda_2\in \K$ we have
$(\lambda_1+\lambda_2)\cdot\uu=\lambda_1\cdot \uu + \lambda_2\cdot\uu$;
\item for any $\uu\in \V$ and any scalars $\lambda_1,\lambda_2\in \K$ we have
$(\lambda_1 \lambda_2)\cdot\uu=\lambda_1\cdot (\lambda_2\cdot\uu)$;
\item for any $\uu\in \V$ we have $1\cdot \uu=\uu$.
\end{itemize}
The elements of $\V$ are called \emph{vectors} and the elements of $\K$ are called \emph{scalars}.
\end{definition}

\begin{fact}[Elementary consequences of the axioms]
Let $\V$ be a vector space over $\K$. For any $\vv\in\V$ and $\lambda\in\K$:
\[
0\cdot \vv = \0, \qquad \lambda\cdot \0 = \0, \qquad (-1)\cdot \vv = -\vv,
\qquad \text{and} \qquad \lambda\cdot\vv = \0 \implies \lambda = 0 \ \text{ or } \ \vv = \0.
\]
\end{fact}

\begin{proof}
We prove the first equality; the others are proved similarly. By distributivity,
\[
0\cdot \vv = (0+0)\cdot \vv = 0\cdot\vv + 0\cdot \vv,
\]
and adding $-(0\cdot\vv)$ to both sides gives $\0 = 0\cdot\vv$.
\end{proof}

\begin{example}[The space $\K^n$]
Let $\K$ be any field and let $n\in \N$. Then
\[
\K^n=\{(x_1,\ldots,x_n)\mid x_i\in \K\}
\]
together with vector addition
\[
(x_1,\ldots,x_n) + (y_1,\ldots,y_n) = (x_1+y_1,\ldots,x_n+y_n)
\]
and scalar multiplication
\[
\lambda\cdot (x_1,\ldots,x_n) = (\lambda x_1,\ldots,\lambda x_n) \qquad (\lambda \in \K)
\]
forms a vector space over $\K$. The zero vector is $\0 = (0,\ldots,0)$.

In the case $\K=\R$ and $n=2$ we recover the well-known space $\R^2$. Taking $\K=\C$ we obtain
$\C^n = \{(z_1,\ldots, z_n)\mid z_i\in\C\}$.
\end{example}

\begin{example}[Bit vectors]
The field $\K$ in the previous example may also be finite. Taking $\K=\Z_2$ gives the space $\Z_2^n$ of
\emph{bit vectors}: vector addition is bitwise XOR, e.g.\ in $\Z_2^4$,
\[
(1,0,1,1)+(1,1,0,1) = (0,1,1,0),
\]
and the only scalars are $0$ and $1$. Note that every vector is its own additive inverse. Identifying a
subset of $\{1,\ldots,n\}$ with its characteristic bit vector, we may equally view $\Z_2^n$ as the set
$2^{\{1,\ldots,n\}}$ of all subsets of $\{1,\ldots,n\}$, with the symmetric difference as vector addition.
Vector spaces over $\Z_2$ are the natural home of error-correcting codes and appear throughout
cryptography.
\end{example}

\begin{example}[Complex numbers as a real vector space]
Recall that the set of complex numbers is $\C = \R\times\R =\{(a,b)\mid a,b\in \R\}$. Together with the
standard addition of complex numbers and multiplication by real scalars, $\C$ can be viewed as a vector
space over the field $\R$.
\end{example}

\begin{example}[Polynomials]
Let $\K$ be a field. The set $\K[x]$ of all polynomials over $\K$, with the standard polynomial addition and
multiplication by a constant from $\K$, forms a vector space over $\K$. The zero vector is the zero
polynomial $0$. Similarly, the set $\K_n[x]$ of polynomials of degree less than or equal to $n$ over $\K$
is a vector space.
\end{example}

\begin{example}[Function spaces]
Let $X$ be a non-empty set and $\K$ a field. The set $\K^X$ of all functions $f\colon X\to \K$, with the
pointwise operations
\[
(f+g)(x) = f(x)+g(x), \qquad (\lambda\cdot f)(x) = \lambda\cdot f(x),
\]
is a vector space over $\K$. The zero vector is the constant function $0$. For $X=\R$ and $\K=\R$ we obtain
the space of all real functions of one real variable.
\end{example}

\section{Subspaces}

\begin{definition}
Let $\V$ be a vector space over $\K$. A subset $\W\subseteq \V$ is called a \emph{subspace} of $\V$ if $\W$ is
itself a vector space over $\K$ under the same operations of vector addition and scalar multiplication.
\end{definition}

\noindent To put it differently, a subset $\W\subseteq \V$ is a subspace of $\V$ if
\begin{itemize}
\item for any $\ww_1,\ww_2\in \W$ we have $\ww_1+\ww_2\in \W$;
\item $\0\in \W$;
\item for any $\ww\in \W$ we have $-\ww\in \W$;
\item for any scalar $\lambda\in \K$ and any $\ww\in \W$ we have $\lambda\cdot \ww\in \W$.
\end{itemize}

\begin{example}
Consider the vector space $\R^3 = \{(x,y,z)\mid x,y,z\in\R\}$. The set
\[
W = \{(x,y,z)\mid 2x+y+3z = 0\}
\]
is a subspace of $\R^3$. Indeed, $\0 = (0,0,0)\in W$; if $\ww_1=(x_1,y_1,z_1)$ and $\ww_2=(x_2,y_2,z_2)$
both satisfy the equation, then so does their sum,
\[
2(x_1+x_2)+(y_1+y_2)+3(z_1+z_2) = \underbrace{(2x_1+y_1+3z_1)}_{=\,0} + \underbrace{(2x_2+y_2+3z_2)}_{=\,0} = 0,
\]
and so does every multiple: $2(\lambda x_1)+\lambda y_1+3(\lambda z_1) = \lambda\cdot 0 = 0$, which covers
both $\lambda\ww_1\in W$ and $-\ww_1\in W$.

On the other hand, the set $W' = \{(x,y,z)\mid x^2 = 1\}$ is \emph{not} a subspace of $\R^3$: it does not
contain the zero vector, and it is not closed under addition either, since $(1,0,0)\in W'$ but
$(1,0,0)+(1,0,0) = (2,0,0)\notin W'$.
\end{example}

\begin{example}
Consider the vector space $\C$ over $\R$. The set $\{z\in\C\mid \Re(z) = 0\}$ is a subspace of $\C$.
\end{example}

\begin{warning}
The intersection of two subspaces of $\V$ is always a subspace, but their \emph{union} usually is not. In
$\R^2$ the two axes $W_1 = \{(x,0)\mid x\in\R\}$ and $W_2 = \{(0,y)\mid y\in\R\}$ are subspaces, yet
$(1,0)+(0,1) = (1,1)\notin W_1\cup W_2$, so $W_1\cup W_2$ is not closed under addition.
\end{warning}

\section{Linear combinations and span}

\begin{definition}
Let $\V$ be a vector space over a field $\K$, let $a_1,\ldots,a_n\in \K$ and $\vv_1,\ldots,\vv_n\in \V$. The vector
\[
a_1\vv_1+\ldots + a_n\vv_n\in \V
\]
is called the \emph{linear combination} of the vectors $\vv_1,\ldots,\vv_n$ with coefficients $a_1,\ldots,a_n$.
\end{definition}

\begin{example}
In $\R^2$ the vector $(2,3)$ is a linear combination of $(1,1)$ and $(0,1)$ with coefficients $2$ and $1$:
\[
(2,3) = 2\cdot (1,1) + 1\cdot (0,1).
\]
\end{example}

\begin{definition}
Let $\V$ be a vector space over a field $\K$ and let $S\subseteq \V$ be a set of vectors. By the \emph{span} of
$S$ we mean the set
\[
\Span{S} = \{a_1\vv_1+\ldots + a_n\vv_n\mid a_i\in \K \text{ and } \vv_i\in S\}\subseteq \V.
\]
In other words, $\Span{S}$ is the set of all possible linear combinations of vectors from $S$.
\end{definition}

\begin{example}
Consider $\R^3$ and the two vectors $(1,2,0),(0,0,3)\in \R^3$. Then
\[
\Span{\{(1,2,0),(0,0,3)\}} = \{ a(1,2,0)+b(0,0,3)\mid a,b\in \R\} = \{ (a,2a,3b)\mid a,b\in \R\}.
\]
\end{example}

\begin{theorem}
Let $\V$ be a vector space over a field $\K$ and let $\vv_1,\ldots,\vv_n\in \V$. Then the set
\[
\Span{\{\vv_1,\ldots,\vv_n\}} = \{a_1\vv_1+\ldots + a_n\vv_n\mid a_i\in \K\}
\]
is a subspace of $\V$ containing the vectors $\vv_1,\ldots,\vv_n$.
\end{theorem}

\begin{proof}
Take $\uu,\ww\in \Span{\{\vv_1,\ldots,\vv_n\}}$, that is
\[
\uu = a_1\vv_1+\ldots + a_n\vv_n, \qquad \ww = b_1\vv_1+\ldots + b_n\vv_n,
\]
for some $a_i,b_i\in \K$. We check the four conditions for a subspace.
\begin{itemize}
\item $\uu+\ww\in \Span{\{\vv_1,\ldots,\vv_n\}}$, because
\[
\uu+\ww = (a_1+b_1)\vv_1+\ldots + (a_n+b_n)\vv_n.
\]
\item $-\uu\in \Span{\{\vv_1,\ldots,\vv_n\}}$, because
\[
-\uu = (-a_1)\vv_1+\ldots + (-a_n)\vv_n.
\]
\item $\0\in \Span{\{\vv_1,\ldots,\vv_n\}}$, because $\0 = 0\cdot\vv_1+\ldots + 0\cdot \vv_n$.
\item $\lambda\cdot\uu\in \Span{\{\vv_1,\ldots,\vv_n\}}$ for any $\lambda\in\K$, because
\[
\lambda\cdot \uu = (\lambda a_1)\vv_1+\ldots + (\lambda a_n)\vv_n. \qedhere
\]
\end{itemize}
\end{proof}

\begin{fact}[Properties of the span]
Let $\V$ be a vector space and let $S,T\subseteq \V$. Then:
\begin{itemize}
\item $S\subseteq \Span{S}$;
\item if $S\subseteq T$, then $\Span{S}\subseteq \Span{T}$;
\item $\Span{\Span{S}} = \Span{S}$;
\item $\Span{S}$ is the \emph{smallest} subspace of $\V$ containing $S$: if $W$ is a subspace of $\V$ and
$S\subseteq W$, then $\Span{S}\subseteq W$.
\end{itemize}
\end{fact}

\begin{example}
\hfill
\begin{itemize}
\item $\C$ over $\R$:
\[
\Span{\{1,i\}} = \{a+bi\mid a,b\in \R\} = \C, \qquad \Span{\{i,2+i\}} = \C,
\]
\[
\Span{\{3+i\}} = \{3a+ai\mid a\in \R\} \neq \C.
\]
\item $\R[x]$ over $\R$:
\[
\Span{\{x^2,x,1\}} = \Span{\{x^2,x,1,x^2+5\}} = \{ax^2+bx+c\mid a,b,c\in \R\},
\]
the space of quadratic functions.
\end{itemize}
\end{example}

\begin{example}[Testing membership in a span]
Is $(1,2,3)\in\Span{\{(1,1,1),(1,0,1)\}}$ in $\R^3$? We look for $a,b\in\R$ with
\[
a(1,1,1)+b(1,0,1) = (1,2,3), \qquad\text{i.e.}\qquad a+b = 1,\quad a = 2,\quad a+b = 3.
\]
The first and third equations contradict each other, so no such $a,b$ exist:
$(1,2,3)\notin \Span{\{(1,1,1),(1,0,1)\}}$. On the other hand, for the vector $(3,1,3)$ the same equations
give $a = 1$, $b = 2$, consistently, hence
\[
(3,1,3) = 1\cdot (1,1,1)+2\cdot (1,0,1) \in \Span{\{(1,1,1),(1,0,1)\}}.
\]
Deciding membership in a span always amounts to solving a system of linear equations; the systematic theory
of such systems is developed in Chapter~\ref{ch:sole}.
\end{example}

\section{Linear independence and basis}

\begin{definition}
Let $\V$ be a vector space over a field $\K$. A set $S\subseteq \V$ of vectors is \emph{linearly independent}
if for any $\vv_1,\ldots,\vv_n \in S$ and any $a_1,\ldots,a_n\in \K$
\[
a_1\vv_1+\ldots+a_n\vv_n = \0 \implies a_1=\ldots =a_n=0.
\]
Otherwise, the set $S$ is \emph{linearly dependent}.
\end{definition}

\begin{example}
In $\R^2$ over $\R$ the set $\{(0,1),(1,1)\}$ is linearly independent, because the linear combination
$a(0,1)+b(1,1) = (b,a+b)$ equals the zero vector $(0,0)$ only if $b=0$ and $a+b=0$, that is $a=b=0$.
\end{example}

\begin{example}
\hfill
\begin{itemize}
\item $\C$ over $\R$: $\{1,i\}$ is linearly independent, since $a\cdot 1 + b\cdot i = 0\implies a=b=0$.
\item $\R[x]$ over $\R$: $\{x+1, x-1, 2\}$ is linearly dependent, because
\[
(x+1)+(-1)(x-1)+(-1)\cdot 2 = 0.
\]
\item $\R[x]$ over $\R$: the set $\{1,x,x^2,x^3,\ldots \}$ is linearly independent.
\end{itemize}
\end{example}

\begin{example}[Testing linear independence]
Are the vectors $(1,2,3)$, $(2,1,0)$, $(1,-4,-9)$ linearly independent in $\R^3$? We solve
\[
a(1,2,3)+b(2,1,0)+c(1,-4,-9) = (0,0,0), \qquad\text{i.e.}\qquad
\begin{aligned}
a+2b+c &= 0,\\ 2a+b-4c &= 0,\\ 3a-9c &= 0.
\end{aligned}
\]
The third equation gives $a = 3c$; substituting into the first gives $b = -2c$, and the second equation is
then satisfied automatically. Taking $c = 1$ produces a non-trivial combination
\[
3\cdot(1,2,3) - 2\cdot (2,1,0) + 1\cdot (1,-4,-9) = (0,0,0),
\]
so the vectors are linearly \emph{dependent}.
\end{example}

\begin{fact}
Let $\V$ be a vector space over $\K$. The set $\{\vv_1,\ldots, \vv_n\}\subseteq \V$ is linearly dependent if and
only if at least one of the vectors $\vv_i$ can be expressed as a linear combination of the others.
\end{fact}

\begin{fact}
Let $\V$ be a vector space over $\K$. If the set $S\subseteq\V$ is linearly independent, then any subset
$S'\subseteq S$ is also linearly independent.
\end{fact}

\begin{definition}
Let $\V$ be a vector space over a field $\K$. A set $B\subseteq \V$ of vectors is called a \emph{basis} of $\V$ if
\begin{itemize}
\item $B$ is linearly independent, and
\item $\Span{B} = \V$.
\end{itemize}
\end{definition}

\begin{example}
\hfill
\begin{itemize}
\item For the vector space $\R^2$ over $\R$ the set $\{(1,0),(1,1)\}$ is a basis.
\item $\C$ over $\R$: the set $\{1,i\}$ is a basis of $\C$; so is, for instance, $\{1,i+2\}$.
\item For the vector space $\K^n$ over $\K$ the set
\[
\{(1,0,\ldots,0),(0,1,\ldots,0),\ldots,(0,0,\ldots,1)\}
\]
is a basis, called the \emph{standard basis} of $\K^n$; its elements are denoted by $e_1,e_2,\ldots,e_n$.
\item $\K_n[x]$ over $\K$: the set $\{1,x,x^2,x^3,\ldots, x^n\}$ is a basis of $\K_n[x]$.
\end{itemize}
\end{example}

\begin{theorem}
Let $\V$ be a vector space over $\K$. A set $\{\vv_1,\ldots,\vv_n\}\subseteq \V$ is a basis of $\V$ if and only if
every vector $\ww\in \V$ can be expressed \emph{uniquely} as a linear combination of $\vv_1,\ldots,\vv_n$.
\end{theorem}

\begin{proof}
($\Rightarrow$) Suppose $\ww=a_1\vv_1+\ldots+a_n\vv_n$ and $\ww=b_1\vv_1+\ldots+b_n\vv_n$. Then
\[
(a_1-b_1)\vv_1+\ldots+(a_n-b_n)\vv_n = \0.
\]
Since $\{\vv_1,\ldots,\vv_n\}$ is linearly independent, $a_1-b_1=0,\ldots,a_n-b_n=0$, that is
$a_1=b_1,\ldots, a_n=b_n$.

($\Leftarrow$) Since every $\ww\in \V$ is a linear combination of $\vv_1,\ldots,\vv_n$, we have
$\Span{\{\vv_1,\ldots,\vv_n\}} =\V$. It remains to show linear independence. The zero vector has the
representation $\0=0\cdot\vv_1+\ldots+0\cdot\vv_n$. Since the representation of every vector is unique,
\[
a_1\vv_1+\ldots+a_n\vv_n =\0 \implies a_1=a_2=\ldots =a_n=0. \qedhere
\]
\end{proof}

\begin{definition}
Let $B = \{\vv_1,\ldots,\vv_n\}$ be a basis of $\V$ (with a fixed order of its elements) and let
$\ww\in \V$. The unique scalars $a_1,\ldots,a_n\in\K$ with $\ww = a_1\vv_1+\ldots+a_n\vv_n$ are called the
\emph{coordinates} of $\ww$ in the basis $B$.
\end{definition}

\begin{example}
In $\R^2$, take the basis $B = \{(1,0),(1,1)\}$ and $\ww = (2,3)$. Solving $(2,3) = a(1,0)+b(1,1)$ gives
$b = 3$ and $a = -1$, so the coordinates of $(2,3)$ in $B$ are $(-1,3)$. In the standard basis
$\{(1,0),(0,1)\}$ the coordinates of $(2,3)$ are simply $(2,3)$.
\end{example}

\section{Dimension}

\begin{theorem}
Let a vector space $\V$ have a finite basis. Then any two bases of $\V$ have the same size.
\end{theorem}

\begin{definition}
If a vector space $\V$ has a finite basis, then the \emph{dimension} of $\V$ is the size of any basis of $\V$.
Otherwise the dimension is defined to be $\infty$. We denote this number by $\dim(\V)$.
\end{definition}

\begin{example}
\hfill
\begin{itemize}
\item $\dim(\{ \0\}) = 0$;
\item $\dim( \K^n) = n$;
\item $\dim(\K[x])=\infty$;
\item $\dim(\K_n[x])=n+1$;
\item $\dim(\C \text{ over } \R)=2$;
\item $\dim(\C \text{ over } \C)=1$.
\end{itemize}
\end{example}

\begin{theorem}
If $\W$ is a subspace of a vector space $\V$, then
\begin{itemize}
\item $\dim(\W) \leq \dim(\V)$;
\item if $\V$ is finite dimensional, then $\dim(\W) = \dim(\V)$ implies $\W=\V$.
\end{itemize}
\end{theorem}

\begin{theorem}
Let $\V$ be a vector space with $\dim(\V)=n$ and let $S\subseteq \V$. Then:
\begin{itemize}
\item if $\Span{S}=\V$, then $|S|\geq n$;
\item if $S$ is linearly independent and $|S|=n$, then $S$ is a basis;
\item if $\Span{S}=\V$ and $|S|=n$, then $S$ is a basis;
\item if $|S|>n$, then $S$ is linearly dependent.
\end{itemize}
\end{theorem}

\begin{theorem}
Let $\V$ be a vector space with $\dim(\V) = n<\infty$. Then:
\begin{itemize}
\item every linearly independent set $S\subseteq \V$ can be extended to a basis of $\V$;
\item every set $S\subseteq \V$ with $\Span{S} = \V$ contains a basis of $\V$.
\end{itemize}
\end{theorem}

\begin{example}
The set $\{(1,1,0)\}$ is linearly independent in $\R^3$; it can be extended to a basis, e.g.\
$\{(1,1,0),(0,1,0),(0,0,1)\}$: a quick check shows that these three vectors are linearly independent, and
by the previous theorem three linearly independent vectors in a $3$-dimensional space form a basis.
\end{example}
